\chapter{Những Câu Chốt Có Duyên}
\chapterintro{Cuối lời phải còn nhịp đẹp}{Một câu chốt gọn và có duyên giúp người dạy giữ được thế mà không làm lớp cụt cảm xúc.}

\section{Thầy cô nói đến đây, phần còn lại để lớp làm cho đẹp}
\begin{manthoai}{Màn thoại}
\textbf{Thầy/Cô:} Thầy cô nói đến đây, phần còn lại để lớp làm cho đẹp.\\
\textbf{Học sinh:} Dạ.
\end{manthoai}
\begin{dungkhinào}
Dùng khi đã nhắc đủ, không muốn nói thêm mà vẫn cần lớp tự chỉnh.
\end{dungkhinào}
\begin{nhipthem}
Câu này hợp để chốt một đoạn nhắc nề nếp mà không kéo thành bài giảng đạo đức.
\end{nhipthem}
\ngancach

\section{Hôm nay lớp làm được, vậy mai đừng làm như chưa từng làm được}
\begin{manthoai}{Màn thoại}
\textbf{Thầy/Cô:} Hôm nay lớp làm được, vậy mai đừng làm như chưa từng làm được.\\
\textbf{Học sinh:} Dạ.
\end{manthoai}
\begin{dungkhinào}
Hợp để chốt sau một buổi lớp phối hợp tốt hoặc sau một lần lớp vừa sửa được lỗi.
\end{dungkhinào}
\begin{nhipthem}
Vừa khen, vừa nhắc, không quá nặng mà vẫn rõ kỳ vọng.
\end{nhipthem}
\ngancach

\section{Thầy cô cần lớp tốt lên, không cần lớp diễn ngoan}
\begin{manthoai}{Màn thoại}
\textbf{Thầy/Cô:} Thầy cô cần lớp tốt lên, không cần lớp diễn ngoan.\\
\textbf{Học sinh:} Dạ.
\end{manthoai}
\begin{dungkhinào}
Rất hợp khi lớp vừa có biểu hiện đối phó, làm cho có lệ.
\end{dungkhinào}
\begin{nhipthem}
Câu này nên dùng khi đã có đủ sự tin cậy với lớp, vì nó đi khá thẳng vào thái độ học sinh.
\end{nhipthem}
\ngancach

\section{Thầy cô nhắc một lần, phần còn lại là uy tín của lớp}
\begin{manthoai}{Màn thoại}
\textbf{Thầy/Cô:} Thầy cô nhắc một lần, phần còn lại là uy tín của lớp.\\
\textbf{Học sinh:} Dạ.
\end{manthoai}
\begin{dungkhinào}
Hợp để chốt sau khi đã nhắc nề nếp đủ rồi và muốn trao phần chủ động lại cho tập thể.
\end{dungkhinào}
\begin{nhipthem}
Câu này gọn, không gay gắt, nhưng tạo cảm giác lớp đang được tin và được giao trách nhiệm.
\end{nhipthem}
\ngancach

\section{Mình khép chuyện ở đây, đừng kéo theo ngày mai nhé}
\begin{manthoai}{Màn thoại}
\textbf{Thầy/Cô:} Mình khép chuyện ở đây, đừng kéo theo ngày mai nhé.\\
\textbf{Học sinh:} Dạ.
\end{manthoai}
\begin{dungkhinào}
Rất hợp khi lớp vừa trải qua một va chạm nhỏ, cần kết lại gọn để không mang tâm trạng sang buổi sau.
\end{dungkhinào}
\begin{nhipthem}
Nên đi cùng một ánh mắt bình tĩnh và một bước chuyển rõ sang hoạt động tiếp theo hoặc lời chào kết.
\end{nhipthem}
